\documentclass[10pt,a4paper]{amsart}
\usepackage{amsmath,amssymb,amsthm,graphicx}
% For commutative diagrams we use Xy-pic
\usepackage[all]{xy}

% We use 2cell for 2-commutative diagrams.
\xyoption{2cell}
% Theorem environments.
%
\theoremstyle{plain}
\newtheorem{theorem}[subsection]{Theorem}
\newtheorem{proposition}[subsection]{Proposition}
\newtheorem{lemma}[subsection]{Lemma}

\theoremstyle{definition}
\newtheorem{definition}[subsection]{Definition}
\newtheorem{example}[subsection]{Example}
\newtheorem{exercise}[subsection]{Exercise}
\newtheorem{situation}[subsection]{Situation}

\theoremstyle{remark}
\newtheorem{remark}[subsection]{Remark}
\newtheorem{remarks}[subsection]{Remarks}

\numberwithin{equation}{subsection}

% Macros
%
\def\lim{\mathop{\mathrm{lim}}\nolimits}
\def\colim{\mathop{\mathrm{colim}}\nolimits}
\def\Spec{\mathop{\mathrm{Spec}}}
\def\Hom{\mathop{\mathrm{Hom}}\nolimits}
\def\Ext{\mathop{\mathrm{Ext}}\nolimits}
\def\SheafHom{\mathop{\mathcal{H}\!\mathit{om}}\nolimits}
\def\SheafExt{\mathop{\mathcal{E}\!\mathit{xt}}\nolimits}
\def\Sch{\mathit{Sch}}
\def\Mor{\mathop{\mathrm{Mor}}\nolimits}
\def\Ob{\mathop{\mathrm{Ob}}\nolimits}
\def\Sh{\mathop{\mathit{Sh}}\nolimits}
\def\NL{\mathop{N\!L}\nolimits}
\def\CH{\mathop{\mathrm{CH}}\nolimits}
\def\proetale{{pro\text{-}\acute{e}tale}}
\def\etale{{\acute{e}tale}}
\def\QCoh{\mathit{QCoh}}
\def\Ker{\mathop{\mathrm{Ker}}}
\def\Im{\mathop{\mathrm{Im}}}
\def\Coker{\mathop{\mathrm{Coker}}}
\def\Coim{\mathop{\mathrm{Coim}}}

% Boxtimes
%
\DeclareMathSymbol{\boxtimes}{\mathbin}{AMSa}{"02}

%
% Macros for moduli stacks/spaces
%
\def\QCohstack{\mathcal{QC}\!\mathit{oh}}
\def\Cohstack{\mathcal{C}\!\mathit{oh}}
\def\Spacesstack{\mathcal{S}\!\mathit{paces}}
\def\Quotfunctor{\mathrm{Quot}}
\def\Hilbfunctor{\mathrm{Hilb}}
\def\Curvesstack{\mathcal{C}\!\mathit{urves}}
\def\Polarizedstack{\mathcal{P}\!\mathit{olarized}}
\def\Complexesstack{\mathcal{C}\!\mathit{omplexes}}
% \Pic is the operator that assigns to X its picard group, usage \Pic(X)
% \Picardstack_{X/B} denotes the Picard stack of X over B
% \Picardfunctor_{X/B} denotes the Picard functor of X over B
\def\Pic{\mathop{\mathrm{Pic}}\nolimits}
\def\Picardstack{\mathcal{P}\!\mathit{ic}}
\def\Picardfunctor{\mathrm{Pic}}
\def\Deformationcategory{\mathcal{D}\!\mathit{ef}}
\newenvironment{reference}{}{}
\newenvironment{slogan}{}{}
\newenvironment{history}{}{}
\usepackage{cleveref}
\crefname{originalsection}{original section}{original sections}
\Crefname{originalsection}{Original section}{Original sections}
\makeatletter
\newlabel{section-de-rham-complex}{{07HX}{source}{}{sourcepointer-1}{}}
\newlabel{section-de-rham-complex@cref}{{[originalsection][0][]07HX}{[1][0][]source}{}{}{}}
\makeatother
\crefname{originalsection}{original section}{original sections}
\Crefname{originalsection}{Original section}{Original sections}
\makeatletter
\newlabel{cohomology-section-cup-product}{{0FKU}{source}{}{sourcepointer-2}{}}
\newlabel{cohomology-section-cup-product@cref}{{[originalsection][0][]0FKU}{[1][0][]source}{}{}{}}
\makeatother
\crefname{originallemma}{original lemma}{original lemmas}
\Crefname{originallemma}{Original lemma}{Original lemmas}
\makeatletter
\newlabel{situation-relative-duality}{{0G8G}{source}{}{sourcepointer-3}{}}
\newlabel{situation-relative-duality@cref}{{[originallemma][0][]0G8G}{[1][0][]source}{}{}{}}
\makeatother
\crefname{originallemma}{original lemma}{original lemmas}
\Crefname{originallemma}{Original lemma}{Original lemmas}
\makeatletter
\newlabel{perfect-lemma-proper-flat-geom-red}{{0G7Y}{source}{}{sourcepointer-4}{}}
\newlabel{perfect-lemma-proper-flat-geom-red@cref}{{[originallemma][0][]0G7Y}{[1][0][]source}{}{}{}}
\makeatother
\crefname{originalsection}{original section}{original sections}
\Crefname{originalsection}{Original section}{Original sections}
\makeatletter
\newlabel{more-morphisms-section-stein-factorization}{{03GX}{source}{}{sourcepointer-5}{}}
\newlabel{more-morphisms-section-stein-factorization@cref}{{[originalsection][0][]03GX}{[1][0][]source}{}{}{}}
\makeatother
\crefname{originallemma}{original lemma}{original lemmas}
\Crefname{originallemma}{Original lemma}{Original lemmas}
\makeatletter
\newlabel{morphisms-lemma-triangle-differentials}{{01UX}{source}{}{sourcepointer-6}{}}
\newlabel{morphisms-lemma-triangle-differentials@cref}{{[originallemma][0][]01UX}{[1][0][]source}{}{}{}}
\makeatother
\crefname{originallemma}{original lemma}{original lemmas}
\Crefname{originallemma}{Original lemma}{Original lemmas}
\makeatletter
\newlabel{morphisms-lemma-etale-at-point}{{02GU}{source}{}{sourcepointer-7}{}}
\newlabel{morphisms-lemma-etale-at-point@cref}{{[originallemma][0][]02GU}{[1][0][]source}{}{}{}}
\makeatother
\crefname{originalremark}{original remark}{original remarks}
\Crefname{originalremark}{Original remark}{Original remarks}
\makeatletter
\newlabel{duality-remark-relative-dualizing-complex}{{0B6S}{source}{}{sourcepointer-8}{}}
\newlabel{duality-remark-relative-dualizing-complex@cref}{{[originalremark][0][]0B6S}{[1][0][]source}{}{}{}}
\makeatother
\crefname{originallemma}{original lemma}{original lemmas}
\Crefname{originallemma}{Original lemma}{Original lemmas}
\makeatletter
\newlabel{duality-lemma-smooth-proper}{{0BRT}{source}{}{sourcepointer-9}{}}
\newlabel{duality-lemma-smooth-proper@cref}{{[originallemma][0][]0BRT}{[1][0][]source}{}{}{}}
\makeatother
\crefname{originallemma}{original lemma}{original lemmas}
\Crefname{originallemma}{Original lemma}{Original lemmas}
\makeatletter
\newlabel{lemma-proper-smooth-de-Rham}{{0FM0}{source}{}{sourcepointer-10}{}}
\newlabel{lemma-proper-smooth-de-Rham@cref}{{[originallemma][0][]0FM0}{[1][0][]source}{}{}{}}
\makeatother
\crefname{originalremark}{original remark}{original remarks}
\Crefname{originalremark}{Original remark}{Original remarks}
\makeatletter
\newlabel{duality-remark-relative-dualizing-complex-relative-cup-product}{{0G81}{source}{}{sourcepointer-11}{}}
\newlabel{duality-remark-relative-dualizing-complex-relative-cup-product@cref}{{[originalremark][0][]0G81}{[1][0][]source}{}{}{}}
\makeatother
\crefname{originallemma}{original lemma}{original lemmas}
\Crefname{originallemma}{Original lemma}{Original lemmas}
\makeatletter
\newlabel{limits-lemma-descend-finite-presentation}{{01ZM}{source}{}{sourcepointer-12}{}}
\newlabel{limits-lemma-descend-finite-presentation@cref}{{[originallemma][0][]01ZM}{[1][0][]source}{}{}{}}
\makeatother
\crefname{originallemma}{original lemma}{original lemmas}
\Crefname{originallemma}{Original lemma}{Original lemmas}
\makeatletter
\newlabel{limits-lemma-descend-smooth}{{0C0C}{source}{}{sourcepointer-13}{}}
\newlabel{limits-lemma-descend-smooth@cref}{{[originallemma][0][]0C0C}{[1][0][]source}{}{}{}}
\makeatother
\crefname{originallemma}{original lemma}{original lemmas}
\Crefname{originallemma}{Original lemma}{Original lemmas}
\makeatletter
\newlabel{limits-lemma-descend-dimension-d}{{0EY2}{source}{}{sourcepointer-14}{}}
\newlabel{limits-lemma-descend-dimension-d@cref}{{[originallemma][0][]0EY2}{[1][0][]source}{}{}{}}
\makeatother
\crefname{originallemma}{original lemma}{original lemmas}
\Crefname{originallemma}{Original lemma}{Original lemmas}
\makeatletter
\newlabel{limits-lemma-descend-surjective}{{07RR}{source}{}{sourcepointer-15}{}}
\newlabel{limits-lemma-descend-surjective@cref}{{[originallemma][0][]07RR}{[1][0][]source}{}{}{}}
\makeatother
\crefname{originallemma}{original lemma}{original lemmas}
\Crefname{originallemma}{Original lemma}{Original lemmas}
\makeatletter
\newlabel{limits-lemma-eventually-proper}{{081F}{source}{}{sourcepointer-16}{}}
\newlabel{limits-lemma-eventually-proper@cref}{{[originallemma][0][]081F}{[1][0][]source}{}{}{}}
\makeatother
\crefname{originallemma}{original lemma}{original lemmas}
\Crefname{originallemma}{Original lemma}{Original lemmas}
\makeatletter
\newlabel{cohomology-lemma-base-change-cup-product}{{0H99}{source}{}{sourcepointer-17}{}}
\newlabel{cohomology-lemma-base-change-cup-product@cref}{{[originallemma][0][]0H99}{[1][0][]source}{}{}{}}
\makeatother
\crefname{originallemma}{original lemma}{original lemmas}
\Crefname{originallemma}{Original lemma}{Original lemmas}
\makeatletter
\newlabel{lemma-top-part-degenerates}{{0FW6}{source}{}{sourcepointer-18}{}}
\newlabel{lemma-top-part-degenerates@cref}{{[originallemma][0][]0FW6}{[1][0][]source}{}{}{}}
\makeatother
\crefname{originallemma}{original lemma}{original lemmas}
\Crefname{originallemma}{Original lemma}{Original lemmas}
\makeatletter
\newlabel{algebra-lemma-reduced-ring-sub-product-fields}{{00EW}{source}{}{sourcepointer-19}{}}
\newlabel{algebra-lemma-reduced-ring-sub-product-fields@cref}{{[originallemma][0][]00EW}{[1][0][]source}{}{}{}}
\makeatother
\crefname{originallemma}{original lemma}{original lemmas}
\Crefname{originallemma}{Original lemma}{Original lemmas}
\makeatletter
\newlabel{algebra-lemma-minimal-prime-reduced-ring}{{00EU}{source}{}{sourcepointer-20}{}}
\newlabel{algebra-lemma-minimal-prime-reduced-ring@cref}{{[originallemma][0][]00EU}{[1][0][]source}{}{}{}}
\makeatother
\crefname{originallemma}{original lemma}{original lemmas}
\Crefname{originallemma}{Original lemma}{Original lemmas}
\makeatletter
\newlabel{coherent-lemma-flat-base-change-cohomology}{{02KH}{source}{}{sourcepointer-21}{}}
\newlabel{coherent-lemma-flat-base-change-cohomology@cref}{{[originallemma][0][]02KH}{[1][0][]source}{}{}{}}
\makeatother
\crefname{originaltheorem}{original theorem}{original theorems}
\Crefname{originaltheorem}{Original theorem}{Original theorems}
\makeatletter
\newlabel{more-morphisms-theorem-stein-factorization-general}{{03H2}{source}{}{sourcepointer-22}{}}
\newlabel{more-morphisms-theorem-stein-factorization-general@cref}{{[originaltheorem][0][]03H2}{[1][0][]source}{}{}{}}
\makeatother
\crefname{originallemma}{original lemma}{original lemmas}
\Crefname{originallemma}{Original lemma}{Original lemmas}
\makeatletter
\newlabel{more-morphisms-lemma-smooth-etale-permanence}{{0G80}{source}{}{sourcepointer-23}{}}
\newlabel{more-morphisms-lemma-smooth-etale-permanence@cref}{{[originallemma][0][]0G80}{[1][0][]source}{}{}{}}
\makeatother
\crefname{originallemma}{original lemma}{original lemmas}
\Crefname{originallemma}{Original lemma}{Original lemmas}
\makeatletter
\newlabel{cohomology-lemma-relative-Leray}{{01F6}{source}{}{sourcepointer-24}{}}
\newlabel{cohomology-lemma-relative-Leray@cref}{{[originallemma][0][]01F6}{[1][0][]source}{}{}{}}
\makeatother
\crefname{originallemma}{original lemma}{original lemmas}
\Crefname{originallemma}{Original lemma}{Original lemmas}
\makeatletter
\newlabel{coherent-lemma-relative-affine-vanishing}{{01XC}{source}{}{sourcepointer-25}{}}
\newlabel{coherent-lemma-relative-affine-vanishing@cref}{{[originallemma][0][]01XC}{[1][0][]source}{}{}{}}
\makeatother
\crefname{originallemma}{original lemma}{original lemmas}
\Crefname{originallemma}{Original lemma}{Original lemmas}
\makeatletter
\newlabel{more-morphisms-lemma-etale-nbhd-dominates-smooth}{{055U}{source}{}{sourcepointer-26}{}}
\newlabel{more-morphisms-lemma-etale-nbhd-dominates-smooth@cref}{{[originallemma][0][]055U}{[1][0][]source}{}{}{}}
\makeatother
\crefname{originalsection}{original section}{original sections}
\Crefname{originalsection}{Original section}{Original sections}
\makeatletter
\newlabel{schemes-lemma-section-immersion}{{01KT}{source}{}{sourcepointer-27}{}}
\newlabel{schemes-lemma-section-immersion@cref}{{[originalsection][0][]01KT}{[1][0][]source}{}{}{}}
\makeatother
\crefname{originallemma}{original lemma}{original lemmas}
\Crefname{originallemma}{Original lemma}{Original lemmas}
\makeatletter
\newlabel{divisors-lemma-regular-quasi-regular-immersion}{{063K}{source}{}{sourcepointer-28}{}}
\newlabel{divisors-lemma-regular-quasi-regular-immersion@cref}{{[originallemma][0][]063K}{[1][0][]source}{}{}{}}
\makeatother
\crefname{originalsection}{original section}{original sections}
\Crefname{originalsection}{Original section}{Original sections}
\makeatletter
\newlabel{divisors-lemma-section-smooth-regular-immersion}{{067R}{source}{}{sourcepointer-29}{}}
\newlabel{divisors-lemma-section-smooth-regular-immersion@cref}{{[originalsection][0][]067R}{[1][0][]source}{}{}{}}
\makeatother
\crefname{originalremark}{original remark}{original remarks}
\Crefname{originalremark}{Original remark}{Original remarks}
\makeatletter
\newlabel{remark-how-to-use}{{0G8A}{source}{}{sourcepointer-30}{}}
\newlabel{remark-how-to-use@cref}{{[originalremark][0][]0G8A}{[1][0][]source}{}{}{}}
\makeatother
\crefname{originallemma}{original lemma}{original lemmas}
\Crefname{originallemma}{Original lemma}{Original lemmas}
\makeatletter
\newlabel{limits-lemma-higher-direct-images-zero-above-dimension-fibre}{{0E7D}{source}{}{sourcepointer-31}{}}
\newlabel{limits-lemma-higher-direct-images-zero-above-dimension-fibre@cref}{{[originallemma][0][]0E7D}{[1][0][]source}{}{}{}}
\makeatother
\crefname{originallemma}{original lemma}{original lemmas}
\Crefname{originallemma}{Original lemma}{Original lemmas}
\makeatletter
\newlabel{lemma-gysin-transverse-global}{{0G8D}{source}{}{sourcepointer-32}{}}
\newlabel{lemma-gysin-transverse-global@cref}{{[originallemma][0][]0G8D}{[1][0][]source}{}{}{}}
\makeatother
\crefname{originallemma}{original lemma}{original lemmas}
\Crefname{originallemma}{Original lemma}{Original lemmas}
\makeatletter
\newlabel{lemma-class-of-a-point}{{0G8E}{source}{}{sourcepointer-33}{}}
\newlabel{lemma-class-of-a-point@cref}{{[originallemma][0][]0G8E}{[1][0][]source}{}{}{}}
\makeatother
\crefname{originallemma}{original lemma}{original lemmas}
\Crefname{originallemma}{Original lemma}{Original lemmas}
\makeatletter
\newlabel{lemma-de-rham-complex-product}{{0FMA}{source}{}{sourcepointer-34}{}}
\newlabel{lemma-de-rham-complex-product@cref}{{[originallemma][0][]0FMA}{[1][0][]source}{}{}{}}
\makeatother
\crefname{originallemma}{original lemma}{original lemmas}
\Crefname{originallemma}{Original lemma}{Original lemmas}
\makeatletter
\newlabel{perfect-lemma-kunneth}{{0FLQ}{source}{}{sourcepointer-35}{}}
\newlabel{perfect-lemma-kunneth@cref}{{[originallemma][0][]0FLQ}{[1][0][]source}{}{}{}}
\makeatother
\crefname{originallemma}{original lemma}{original lemmas}
\Crefname{originallemma}{Original lemma}{Original lemmas}
\makeatletter
\newlabel{perfect-lemma-kunneth-single-sheaf}{{0FU4}{source}{}{sourcepointer-36}{}}
\newlabel{perfect-lemma-kunneth-single-sheaf@cref}{{[originallemma][0][]0FU4}{[1][0][]source}{}{}{}}
\makeatother
\crefname{originallemma}{original lemma}{original lemmas}
\Crefname{originallemma}{Original lemma}{Original lemmas}
\makeatletter
\newlabel{lemma-gysin-differential-hodge}{{0G8B}{source}{}{sourcepointer-37}{}}
\newlabel{lemma-gysin-differential-hodge@cref}{{[originallemma][0][]0G8B}{[1][0][]source}{}{}{}}
\makeatother
\crefname{originalsection}{original section}{original sections}
\Crefname{originalsection}{Original section}{Original sections}
\makeatletter
\newlabel{section-kunneth}{{0FM9}{source}{}{sourcepointer-38}{}}
\newlabel{section-kunneth@cref}{{[originalsection][0][]0FM9}{[1][0][]source}{}{}{}}
\makeatother
\crefname{originalsection}{original section}{original sections}
\Crefname{originalsection}{Original section}{Original sections}
\makeatletter
\newlabel{perfect-section-kunneth-complexes}{{0G4A}{source}{}{sourcepointer-39}{}}
\newlabel{perfect-section-kunneth-complexes@cref}{{[originalsection][0][]0G4A}{[1][0][]source}{}{}{}}
\makeatother
\crefname{originallemma}{original lemma}{original lemmas}
\Crefname{originallemma}{Original lemma}{Original lemmas}
\makeatletter
\newlabel{lemma-gysin-projection-global}{{0G8C}{source}{}{sourcepointer-40}{}}
\newlabel{lemma-gysin-projection-global@cref}{{[originallemma][0][]0G8C}{[1][0][]source}{}{}{}}
\makeatother
\crefname{originalremark}{original remark}{original remarks}
\Crefname{originalremark}{Original remark}{Original remarks}
\makeatletter
\newlabel{cohomology-remark-cup-with-element-map-total-cohomology}{{0G6W}{source}{}{sourcepointer-41}{}}
\newlabel{cohomology-remark-cup-with-element-map-total-cohomology@cref}{{[originalremark][0][]0G6W}{[1][0][]source}{}{}{}}
\makeatother
\crefname{originalremark}{original remark}{original remarks}
\Crefname{originalremark}{Original remark}{Original remarks}
\makeatletter
\newlabel{perfect-remark-annoying-compatibility}{{0G7W}{source}{}{sourcepointer-42}{}}
\newlabel{perfect-remark-annoying-compatibility@cref}{{[originalremark][0][]0G7W}{[1][0][]source}{}{}{}}
\makeatother
\crefname{originallemma}{original lemma}{original lemmas}
\Crefname{originallemma}{Original lemma}{Original lemmas}
\makeatletter
\newlabel{more-algebra-lemma-splitting-unique}{{0BCG}{source}{}{sourcepointer-43}{}}
\newlabel{more-algebra-lemma-splitting-unique@cref}{{[originallemma][0][]0BCG}{[1][0][]source}{}{}{}}
\makeatother
\crefname{originalproposition}{original proposition}{original propositions}
\Crefname{originalproposition}{Original proposition}{Original propositions}
\makeatletter
\newlabel{proposition-poincare-duality}{{0FW7}{source}{}{sourcepointer-44}{}}
\newlabel{proposition-poincare-duality@cref}{{[originalproposition][0][]0FW7}{[1][0][]source}{}{}{}}
\makeatother
\crefname{originalsection}{original section}{original sections}
\Crefname{originalsection}{Original section}{Original sections}
\makeatletter
\newlabel{section-hodge-to-de-rham}{{0FM6}{source}{}{sourcepointer-45}{}}
\newlabel{section-hodge-to-de-rham@cref}{{[originalsection][0][]0FM6}{[1][0][]source}{}{}{}}
\makeatother
\crefname{originallemma}{original lemma}{original lemmas}
\Crefname{originallemma}{Original lemma}{Original lemmas}
\makeatletter
\newlabel{lemma-duality-hodge}{{0FW4}{source}{}{sourcepointer-46}{}}
\newlabel{lemma-duality-hodge@cref}{{[originallemma][0][]0FW4}{[1][0][]source}{}{}{}}
\makeatother
\crefname{originallemma}{original lemma}{original lemmas}
\Crefname{originallemma}{Original lemma}{Original lemmas}
\makeatletter
\newlabel{derived-lemma-third-isomorphism-triangle}{{014A}{source}{}{sourcepointer-47}{}}
\newlabel{derived-lemma-third-isomorphism-triangle@cref}{{[originallemma][0][]014A}{[1][0][]source}{}{}{}}
\makeatother
\title{Relative de Rham Poincare duality as a perfect pairing of perfect complexes}
\author{The Stacks Project Authors}
\date{stacks-project:derham.tex:local-snapshot-20261011; GNU FDL 1.2 or later}
\begin{document}
\maketitle
\noindent\textit{Editorial note.} Original author theorem and proof are retained. Only this target is counted; local supporting statements are not separate new entries.\par\medskip
\noindent\textit{Allowed standard prerequisites (editorial).}\begin{itemize}
\item Derived categories, perfect complexes and duality, relative Hodge-to-de Rham spectral sequences, relative coherent cohomological dimension and Serre duality. The local relative Hodge and boundary-degeneration arguments are retained.. These are the named standard inputs; the selected main construction and its local nonroutine support remain in full.
\end{itemize}\medskip
\ifdefined\hypertarget\hypertarget{sourcepointer-1}{}\fi
\noindent\textit{Source reference.} Original section 07HX: focused-seeds/AGP/sources/stacks-snapshot-20261011/tags.txt, source lines 9567--9567. Stable Stacks locator for declared standard input or surrounding context. No selected local proof is replaced by this pointer.\par\medskip
\ifdefined\hypertarget\hypertarget{sourcepointer-2}{}\fi
\noindent\textit{Source reference.} Original section 0FKU: focused-seeds/AGP/sources/stacks-snapshot-20261011/tags.txt, source lines 19474--19474. Stable Stacks locator for declared standard input or surrounding context. No selected local proof is replaced by this pointer.\par\medskip
\ifdefined\hypertarget\hypertarget{sourcepointer-3}{}\fi
\noindent\textit{Source reference.} Original lemma 0G8G: focused-seeds/AGP/sources/stacks-snapshot-20261011/tags.txt, source lines 20266--20266. Stable Stacks locator for declared standard input or surrounding context. No selected local proof is replaced by this pointer.\par\medskip
\ifdefined\hypertarget\hypertarget{sourcepointer-4}{}\fi
\noindent\textit{Source reference.} Original lemma 0G7Y: focused-seeds/AGP/sources/stacks-snapshot-20261011/tags.txt, source lines 20248--20248. Stable Stacks locator for declared standard input or surrounding context. No selected local proof is replaced by this pointer.\par\medskip
\ifdefined\hypertarget\hypertarget{sourcepointer-5}{}\fi
\noindent\textit{Source reference.} Original section 03GX: focused-seeds/AGP/sources/stacks-snapshot-20261011/tags.txt, source lines 4544--4544. Stable Stacks locator for declared standard input or surrounding context. No selected local proof is replaced by this pointer.\par\medskip
\ifdefined\hypertarget\hypertarget{sourcepointer-6}{}\fi
\noindent\textit{Source reference.} Original lemma 01UX: focused-seeds/AGP/sources/stacks-snapshot-20261011/tags.txt, source lines 2426--2426. Stable Stacks locator for declared standard input or surrounding context. No selected local proof is replaced by this pointer.\par\medskip
\ifdefined\hypertarget\hypertarget{sourcepointer-7}{}\fi
\noindent\textit{Source reference.} Original lemma 02GU: focused-seeds/AGP/sources/stacks-snapshot-20261011/tags.txt, source lines 3241--3241. Stable Stacks locator for declared standard input or surrounding context. No selected local proof is replaced by this pointer.\par\medskip
\ifdefined\hypertarget\hypertarget{sourcepointer-8}{}\fi
\noindent\textit{Source reference.} Original remark 0B6S: focused-seeds/AGP/sources/stacks-snapshot-20261011/tags.txt, source lines 14080--14080. Stable Stacks locator for declared standard input or surrounding context. No selected local proof is replaced by this pointer.\par\medskip
\ifdefined\hypertarget\hypertarget{sourcepointer-9}{}\fi
\noindent\textit{Source reference.} Original lemma 0BRT: focused-seeds/AGP/sources/stacks-snapshot-20261011/tags.txt, source lines 14781--14781. Stable Stacks locator for declared standard input or surrounding context. No selected local proof is replaced by this pointer.\par\medskip
\ifdefined\hypertarget\hypertarget{sourcepointer-10}{}\fi
\noindent\textit{Source reference.} Original lemma 0FM0: focused-seeds/AGP/sources/stacks-snapshot-20261011/tags.txt, source lines 19515--19515. Stable Stacks locator for declared standard input or surrounding context. No selected local proof is replaced by this pointer.\par\medskip
\ifdefined\hypertarget\hypertarget{sourcepointer-11}{}\fi
\noindent\textit{Source reference.} Original remark 0G81: focused-seeds/AGP/sources/stacks-snapshot-20261011/tags.txt, source lines 20251--20251. Stable Stacks locator for declared standard input or surrounding context. No selected local proof is replaced by this pointer.\par\medskip
\ifdefined\hypertarget\hypertarget{sourcepointer-12}{}\fi
\noindent\textit{Source reference.} Original lemma 01ZM: focused-seeds/AGP/sources/stacks-snapshot-20261011/tags.txt, source lines 2599--2599. Stable Stacks locator for declared standard input or surrounding context. No selected local proof is replaced by this pointer.\par\medskip
\ifdefined\hypertarget\hypertarget{sourcepointer-13}{}\fi
\noindent\textit{Source reference.} Original lemma 0C0C: focused-seeds/AGP/sources/stacks-snapshot-20261011/tags.txt, source lines 15082--15082. Stable Stacks locator for declared standard input or surrounding context. No selected local proof is replaced by this pointer.\par\medskip
\ifdefined\hypertarget\hypertarget{sourcepointer-14}{}\fi
\noindent\textit{Source reference.} Original lemma 0EY2: focused-seeds/AGP/sources/stacks-snapshot-20261011/tags.txt, source lines 18677--18677. Stable Stacks locator for declared standard input or surrounding context. No selected local proof is replaced by this pointer.\par\medskip
\ifdefined\hypertarget\hypertarget{sourcepointer-15}{}\fi
\noindent\textit{Source reference.} Original lemma 07RR: focused-seeds/AGP/sources/stacks-snapshot-20261011/tags.txt, source lines 9876--9876. Stable Stacks locator for declared standard input or surrounding context. No selected local proof is replaced by this pointer.\par\medskip
\ifdefined\hypertarget\hypertarget{sourcepointer-16}{}\fi
\noindent\textit{Source reference.} Original lemma 081F: focused-seeds/AGP/sources/stacks-snapshot-20261011/tags.txt, source lines 10215--10215. Stable Stacks locator for declared standard input or surrounding context. No selected local proof is replaced by this pointer.\par\medskip
\ifdefined\hypertarget\hypertarget{sourcepointer-17}{}\fi
\noindent\textit{Source reference.} Original lemma 0H99: focused-seeds/AGP/sources/stacks-snapshot-20261011/tags.txt, source lines 21519--21519. Stable Stacks locator for declared standard input or surrounding context. No selected local proof is replaced by this pointer.\par\medskip
\ifdefined\hypertarget\hypertarget{sourcepointer-18}{}\fi
\noindent\textit{Source reference.} Original lemma 0FW6: focused-seeds/AGP/sources/stacks-snapshot-20261011/tags.txt, source lines 19836--19836. Stable Stacks locator for declared standard input or surrounding context. No selected local proof is replaced by this pointer.\par\medskip
\ifdefined\hypertarget\hypertarget{sourcepointer-19}{}\fi
\noindent\textit{Source reference.} Original lemma 00EW: focused-seeds/AGP/sources/stacks-snapshot-20261011/tags.txt, source lines 550--550. Stable Stacks locator for declared standard input or surrounding context. No selected local proof is replaced by this pointer.\par\medskip
\ifdefined\hypertarget\hypertarget{sourcepointer-20}{}\fi
\noindent\textit{Source reference.} Original lemma 00EU: focused-seeds/AGP/sources/stacks-snapshot-20261011/tags.txt, source lines 548--548. Stable Stacks locator for declared standard input or surrounding context. No selected local proof is replaced by this pointer.\par\medskip
\ifdefined\hypertarget\hypertarget{sourcepointer-21}{}\fi
\noindent\textit{Source reference.} Original lemma 02KH: focused-seeds/AGP/sources/stacks-snapshot-20261011/tags.txt, source lines 3372--3372. Stable Stacks locator for declared standard input or surrounding context. No selected local proof is replaced by this pointer.\par\medskip
\ifdefined\hypertarget\hypertarget{sourcepointer-22}{}\fi
\noindent\textit{Source reference.} Original theorem 03H2: focused-seeds/AGP/sources/stacks-snapshot-20261011/tags.txt, source lines 4549--4549. Stable Stacks locator for declared standard input or surrounding context. No selected local proof is replaced by this pointer.\par\medskip
\ifdefined\hypertarget\hypertarget{sourcepointer-23}{}\fi
\noindent\textit{Source reference.} Original lemma 0G80: focused-seeds/AGP/sources/stacks-snapshot-20261011/tags.txt, source lines 20250--20250. Stable Stacks locator for declared standard input or surrounding context. No selected local proof is replaced by this pointer.\par\medskip
\ifdefined\hypertarget\hypertarget{sourcepointer-24}{}\fi
\noindent\textit{Source reference.} Original lemma 01F6: focused-seeds/AGP/sources/stacks-snapshot-20261011/tags.txt, source lines 1859--1859. Stable Stacks locator for declared standard input or surrounding context. No selected local proof is replaced by this pointer.\par\medskip
\ifdefined\hypertarget\hypertarget{sourcepointer-25}{}\fi
\noindent\textit{Source reference.} Original lemma 01XC: focused-seeds/AGP/sources/stacks-snapshot-20261011/tags.txt, source lines 2513--2513. Stable Stacks locator for declared standard input or surrounding context. No selected local proof is replaced by this pointer.\par\medskip
\ifdefined\hypertarget\hypertarget{sourcepointer-26}{}\fi
\noindent\textit{Source reference.} Original lemma 055U: focused-seeds/AGP/sources/stacks-snapshot-20261011/tags.txt, source lines 6687--6687. Stable Stacks locator for declared standard input or surrounding context. No selected local proof is replaced by this pointer.\par\medskip
\ifdefined\hypertarget\hypertarget{sourcepointer-27}{}\fi
\noindent\textit{Source reference.} Original section 01KT: focused-seeds/AGP/sources/stacks-snapshot-20261011/tags.txt, source lines 2062--2062. Stable Stacks locator for declared standard input or surrounding context. No selected local proof is replaced by this pointer.\par\medskip
\ifdefined\hypertarget\hypertarget{sourcepointer-28}{}\fi
\noindent\textit{Source reference.} Original lemma 063K: focused-seeds/AGP/sources/stacks-snapshot-20261011/tags.txt, source lines 7834--7834. Stable Stacks locator for declared standard input or surrounding context. No selected local proof is replaced by this pointer.\par\medskip
\ifdefined\hypertarget\hypertarget{sourcepointer-29}{}\fi
\noindent\textit{Source reference.} Original section 067R: focused-seeds/AGP/sources/stacks-snapshot-20261011/tags.txt, source lines 7980--7980. Stable Stacks locator for declared standard input or surrounding context. No selected local proof is replaced by this pointer.\par\medskip
\ifdefined\hypertarget\hypertarget{sourcepointer-30}{}\fi
\noindent\textit{Source reference.} Original remark 0G8A: focused-seeds/AGP/sources/stacks-snapshot-20261011/tags.txt, source lines 20260--20260. Stable Stacks locator for declared standard input or surrounding context. No selected local proof is replaced by this pointer.\par\medskip
\ifdefined\hypertarget\hypertarget{sourcepointer-31}{}\fi
\noindent\textit{Source reference.} Original lemma 0E7D: focused-seeds/AGP/sources/stacks-snapshot-20261011/tags.txt, source lines 17778--17778. Stable Stacks locator for declared standard input or surrounding context. No selected local proof is replaced by this pointer.\par\medskip
\ifdefined\hypertarget\hypertarget{sourcepointer-32}{}\fi
\noindent\textit{Source reference.} Original lemma 0G8D: focused-seeds/AGP/sources/stacks-snapshot-20261011/tags.txt, source lines 20263--20263. Stable Stacks locator for declared standard input or surrounding context. No selected local proof is replaced by this pointer.\par\medskip
\ifdefined\hypertarget\hypertarget{sourcepointer-33}{}\fi
\noindent\textit{Source reference.} Original lemma 0G8E: focused-seeds/AGP/sources/stacks-snapshot-20261011/tags.txt, source lines 20264--20264. Stable Stacks locator for declared standard input or surrounding context. No selected local proof is replaced by this pointer.\par\medskip
\ifdefined\hypertarget\hypertarget{sourcepointer-34}{}\fi
\noindent\textit{Source reference.} Original lemma 0FMA: focused-seeds/AGP/sources/stacks-snapshot-20261011/tags.txt, source lines 19525--19525. Stable Stacks locator for declared standard input or surrounding context. No selected local proof is replaced by this pointer.\par\medskip
\ifdefined\hypertarget\hypertarget{sourcepointer-35}{}\fi
\noindent\textit{Source reference.} Original lemma 0FLQ: focused-seeds/AGP/sources/stacks-snapshot-20261011/tags.txt, source lines 19505--19505. Stable Stacks locator for declared standard input or surrounding context. No selected local proof is replaced by this pointer.\par\medskip
\ifdefined\hypertarget\hypertarget{sourcepointer-36}{}\fi
\noindent\textit{Source reference.} Original lemma 0FU4: focused-seeds/AGP/sources/stacks-snapshot-20261011/tags.txt, source lines 19764--19764. Stable Stacks locator for declared standard input or surrounding context. No selected local proof is replaced by this pointer.\par\medskip
\ifdefined\hypertarget\hypertarget{sourcepointer-37}{}\fi
\noindent\textit{Source reference.} Original lemma 0G8B: focused-seeds/AGP/sources/stacks-snapshot-20261011/tags.txt, source lines 20261--20261. Stable Stacks locator for declared standard input or surrounding context. No selected local proof is replaced by this pointer.\par\medskip
\ifdefined\hypertarget\hypertarget{sourcepointer-38}{}\fi
\noindent\textit{Source reference.} Original section 0FM9: focused-seeds/AGP/sources/stacks-snapshot-20261011/tags.txt, source lines 19524--19524. Stable Stacks locator for declared standard input or surrounding context. No selected local proof is replaced by this pointer.\par\medskip
\ifdefined\hypertarget\hypertarget{sourcepointer-39}{}\fi
\noindent\textit{Source reference.} Original section 0G4A: focused-seeds/AGP/sources/stacks-snapshot-20261011/tags.txt, source lines 20120--20120. Stable Stacks locator for declared standard input or surrounding context. No selected local proof is replaced by this pointer.\par\medskip
\ifdefined\hypertarget\hypertarget{sourcepointer-40}{}\fi
\noindent\textit{Source reference.} Original lemma 0G8C: focused-seeds/AGP/sources/stacks-snapshot-20261011/tags.txt, source lines 20262--20262. Stable Stacks locator for declared standard input or surrounding context. No selected local proof is replaced by this pointer.\par\medskip
\ifdefined\hypertarget\hypertarget{sourcepointer-41}{}\fi
\noindent\textit{Source reference.} Original remark 0G6W: focused-seeds/AGP/sources/stacks-snapshot-20261011/tags.txt, source lines 20211--20211. Stable Stacks locator for declared standard input or surrounding context. No selected local proof is replaced by this pointer.\par\medskip
\ifdefined\hypertarget\hypertarget{sourcepointer-42}{}\fi
\noindent\textit{Source reference.} Original remark 0G7W: focused-seeds/AGP/sources/stacks-snapshot-20261011/tags.txt, source lines 20246--20246. Stable Stacks locator for declared standard input or surrounding context. No selected local proof is replaced by this pointer.\par\medskip
\ifdefined\hypertarget\hypertarget{sourcepointer-43}{}\fi
\noindent\textit{Source reference.} Original lemma 0BCG: focused-seeds/AGP/sources/stacks-snapshot-20261011/tags.txt, source lines 14279--14279. Stable Stacks locator for declared standard input or surrounding context. No selected local proof is replaced by this pointer.\par\medskip
\ifdefined\hypertarget\hypertarget{sourcepointer-44}{}\fi
\noindent\textit{Source reference.} Original proposition 0FW7: focused-seeds/AGP/sources/stacks-snapshot-20261011/tags.txt, source lines 19837--19837. Stable Stacks locator for declared standard input or surrounding context. No selected local proof is replaced by this pointer.\par\medskip
\ifdefined\hypertarget\hypertarget{sourcepointer-45}{}\fi
\noindent\textit{Source reference.} Original section 0FM6: focused-seeds/AGP/sources/stacks-snapshot-20261011/tags.txt, source lines 19521--19521. Stable Stacks locator for declared standard input or surrounding context. No selected local proof is replaced by this pointer.\par\medskip
\ifdefined\hypertarget\hypertarget{sourcepointer-46}{}\fi
\noindent\textit{Source reference.} Original lemma 0FW4: focused-seeds/AGP/sources/stacks-snapshot-20261011/tags.txt, source lines 19834--19834. Stable Stacks locator for declared standard input or surrounding context. No selected local proof is replaced by this pointer.\par\medskip
\ifdefined\hypertarget\hypertarget{sourcepointer-47}{}\fi
\noindent\textit{Source reference.} Original lemma 014A: focused-seeds/AGP/sources/stacks-snapshot-20261011/tags.txt, source lines 1465--1465. Stable Stacks locator for declared standard input or surrounding context. No selected local proof is replaced by this pointer.\par\medskip
\section{Cup product}
\label{section-cup-product}

\noindent
Consider the maps
$\Omega^p_{X/S} \times \Omega^q_{X/S} \to \Omega^{p + q}_{X/S}$
given by $(\omega , \eta) \longmapsto \omega \wedge \eta$.
Using the formula for $\text{d}$ given in Section \ref{section-de-rham-complex}
and the Leibniz rule for $\text{d} : \mathcal{O}_X \to \Omega_{X/S}$
we see that $\text{d}(\omega \wedge \eta) = \text{d}(\omega) \wedge \eta +
(-1)^{\deg(\omega)} \omega \wedge \text{d}(\eta)$. This means that
$\wedge$ defines a morphism
\begin{equation}
\label{equation-wedge}
\wedge :
\text{Tot}(
\Omega^\bullet_{X/S} \otimes_{p^{-1}\mathcal{O}_S} \Omega^\bullet_{X/S})
\longrightarrow
\Omega^\bullet_{X/S}
\end{equation}
of complexes of $p^{-1}\mathcal{O}_S$-modules.

\medskip\noindent
Combining the cup product of
Cohomology, Section \ref{cohomology-section-cup-product}
with (\ref{equation-wedge}) we find a
$H^0(S, \mathcal{O}_S)$-bilinear cup product map
$$
\cup : H^i_{dR}(X/S) \times H^j_{dR}(X/S) \longrightarrow H^{i + j}_{dR}(X/S)
$$
For example, if $\omega \in \Gamma(X, \Omega^i_{X/S})$ and
$\eta \in \Gamma(X, \Omega^j_{X/S})$ are closed, then
the cup product of the de Rham cohomology classes of
$\omega$ and $\eta$ is the de Rham cohomology class of $\omega \wedge \eta$,
see discussion in Cohomology, Section \ref{cohomology-section-cup-product}.

\medskip\noindent
Given a commutative diagram
$$
\xymatrix{
X' \ar[r]_f \ar[d] & X \ar[d] \\
S' \ar[r] & S
}
$$
of schemes, the pullback maps
$f^* : R\Gamma(X, \Omega^\bullet_{X/S}) \to R\Gamma(X', \Omega^\bullet_{X'/S'})$
and
$f^* : H^i_{dR}(X/S) \longrightarrow H^i_{dR}(X'/S')$
are compatible with the cup product defined above.

\begin{lemma}
\label{lemma-relative-bottom-part-degenerates}
In Situation \ref{situation-relative-duality} the pushforward
$f_*\mathcal{O}_X$ is a finite \'etale $\mathcal{O}_S$-algebra
and locally on $S$ we have $Rf_*\mathcal{O}_X = f_*\mathcal{O}_X \oplus P$
in $D(\mathcal{O}_S)$ with $P$ perfect of tor amplitude in $[1, \infty)$.
The map $\text{d} : f_*\mathcal{O}_X \to f_*\Omega_{X/S}$ is zero.
\end{lemma}

\begin{proof}
The first part of the statement follows from
Derived Categories of Schemes, Lemma \ref{perfect-lemma-proper-flat-geom-red}.
Setting $S' = \underline{\Spec}_S(f_*\mathcal{O}_X)$ we get a factorization
$X \to S' \to S$ (this is the Stein factorization, see
More on Morphisms, Section \ref{more-morphisms-section-stein-factorization},
although we don't need this)
and we see that $\Omega_{X/S} = \Omega_{X/S'}$ for example by
Morphisms, Lemma \ref{morphisms-lemma-triangle-differentials} and
\ref{morphisms-lemma-etale-at-point}. This of course implies that
$\text{d} : f_*\mathcal{O}_X \to f_*\Omega_{X/S}$ is zero.
\end{proof}

\begin{lemma}
\label{lemma-relative-duality-hodge}
In Situation \ref{situation-relative-duality} there exists an
$\mathcal{O}_S$-module map
$$
t : Rf_*\Omega^n_{X/S}[n] \longrightarrow \mathcal{O}_S
$$
unique up to precomposing by multiplication by a unit of
$H^0(X, \mathcal{O}_X)$ with the following property: for all $p$ the pairing
$$
Rf_*\Omega^p_{X/S}
\otimes_{\mathcal{O}_S}^\mathbf{L}
Rf_*\Omega^{n - p}_{X/S}[n]
\longrightarrow
\mathcal{O}_S
$$
given by the relative cup product composed with $t$
is a perfect pairing of perfect complexes on $S$.
\end{lemma}

\begin{proof}
We first assume $S$ is a Noetherian scheme.
Let $\omega^\bullet_{X/S}$ be the relative dualizing complex of $X$ over $S$ as
in Duality for Schemes, Remark \ref{duality-remark-relative-dualizing-complex}
and let $Rf_*\omega_{X/S}^\bullet \to \mathcal{O}_S$ be its trace map. By
Duality for Schemes, Lemma \ref{duality-lemma-smooth-proper}
(here we use that $S$ is Noetherian)
there exists an isomorphism $\omega^\bullet_{X/S} \cong \Omega^n_{X/S}[n]$
and using this isomorphism we obtain $t$. The complexes $Rf_*\Omega^p_{X/S}$
are perfect by Lemma \ref{lemma-proper-smooth-de-Rham}.
Since $\Omega^p_{X/S}$ is locally free and since
$\Omega^p_{X/S} \otimes_{\mathcal{O}_X} \Omega^{n - p}_{X/S} \to
\Omega^n_{X/S}$ exhibits an isomorphism $\Omega^p_{X/S} \cong
\SheafHom_{\mathcal{O}_X}(\Omega^{n - p}_{X/S}, \Omega^n_{X/S})$
we see that the pairing induced by the relative cup product is perfect by
Duality for Schemes, Remark
\ref{duality-remark-relative-dualizing-complex-relative-cup-product}.

\medskip\noindent
Proof of existence in the general case.
By absolute Noetherian approximation we can find a cartesian diagram
$$
\xymatrix{
X \ar[r] \ar[d]^f & X_0 \ar[d]^{f_0} \\
S \ar[r]^g & S_0
}
$$
with $S_0$ Noetherian and $f_0$ a proper and smooth morphism
all of whose fibres are nonempty and equidimensional of dimension $n$.
See Limits, Lemmas \ref{limits-lemma-descend-finite-presentation},
\ref{limits-lemma-descend-smooth}, \ref{limits-lemma-descend-dimension-d},
\ref{limits-lemma-descend-surjective}, and
\ref{limits-lemma-eventually-proper}. Choose a map
$$
t_0 : Rf_{0, *}\Omega^n_{X_0/S_0}[n] \longrightarrow \mathcal{O}_{S_0}
$$
as in the statement of the lemma. By cohomology and base change the
complexes $Rf_*(\Omega^p_{X/S})$ are equal to
$Lg^*(Rf_{0, *}\Omega^p_{X_0/S_0})$ for all $p \geq 0$, see
Lemma \ref{lemma-proper-smooth-de-Rham}. In particular the pullback
of $t_0$ by $g$ gives us a map $t$ as in the statement of the lemma.
Namely, by Cohomology, Lemma \ref{cohomology-lemma-base-change-cup-product},
the cup product map
$$
Rf_*\Omega^p_{X/S}
\otimes_{\mathcal{O}_S}^\mathbf{L}
Rf_*\Omega^{n - p}_{X/S}[n]
\longrightarrow
Rf_*\Omega^n_{X/S}
$$
is the pullback by $g$ of the corresponding cup product map for $f_0$.
Hence the fact that $t_0$ produces perfect pairings, pulls back to
the same property for $t$ on $S$.

\medskip\noindent
Uniqueness of $t$. Choose a distinguished triangle
$f_*\mathcal{O}_X \to Rf_*\mathcal{O}_X \to P \to f_*\mathcal{O}_X[1]$.
By Lemma \ref{lemma-relative-bottom-part-degenerates}
the object $P$ is perfect of tor amplitude in $[1, \infty)$
and the triangle is locally on $S$ split.
Thus $R\SheafHom_{\mathcal{O}_X}(P, \mathcal{O}_X)$ is perfect
of tor amplitude in $(-\infty, -1]$. Hence duality (above) shows that
locally on $S$ we have
$$
Rf_*\Omega^n_{X/S}[n] \cong
R\SheafHom_{\mathcal{O}_S}(f_*\mathcal{O}_X, \mathcal{O}_S)
\oplus R\SheafHom_{\mathcal{O}_X}(P, \mathcal{O}_X)
$$
This shows that $R^nf_*\Omega^n_{X/S}$ is finite locally free and
that we obtain a perfect $\mathcal{O}_S$-bilinear pairing
$$
f_*\mathcal{O}_X \times R^nf_*\Omega^n_{X/S} \longrightarrow \mathcal{O}_S
$$
using $t$.
This implies that any $\mathcal{O}_S$-linear map
$t' : R^nf_*\Omega^n_{X/S} \to \mathcal{O}_S$ is of the form
$t' = t \circ g$ for some
$g \in \Gamma(S, f_*\mathcal{O}_X) = \Gamma(X, \mathcal{O}_X)$.
In order for $t'$ to still determine a perfect pairing $g$ will have
to be a unit. This finishes the proof.
\end{proof}

\begin{lemma}
\label{lemma-relative-top-part-degenerates}
In Situation \ref{situation-relative-duality} the map
$\text{d} : R^nf_*\Omega^{n - 1}_{X/S} \to R^nf_*\Omega^n_{X/S}$
is zero.
\end{lemma}

\noindent
As we mentioned in the proof of Lemma \ref{lemma-top-part-degenerates}
this lemma is not an easy consequence of Lemmas
\ref{lemma-relative-duality-hodge} and
\ref{lemma-relative-bottom-part-degenerates}.

\begin{proof}[Proof in case $S$ is reduced]
Assume $S$ is reduced. Observe that
$\text{d} : R^nf_*\Omega^{n - 1}_{X/S} \to R^nf_*\Omega^n_{X/S}$
is an $\mathcal{O}_S$-linear map of (quasi-coherent) $\mathcal{O}_S$-modules.
The $\mathcal{O}_S$-module $R^nf_*\Omega^n_{X/S}$ is finite locally free
(as the dual of the finite locally free $\mathcal{O}_S$-module
$f_*\mathcal{O}_X$ by Lemmas
\ref{lemma-relative-duality-hodge} and
\ref{lemma-relative-bottom-part-degenerates}).
Since $S$ is reduced it suffices to show that
the stalk of $\text{d}$ in every generic point $\eta \in S$
is zero; this follows by looking at sections over affine opens,
using that the target of $\text{d}$ is locally free, and
Algebra, Lemma \ref{algebra-lemma-reduced-ring-sub-product-fields} part (2).
Since $S$ is reduced we have $\mathcal{O}_{S, \eta} = \kappa(\eta)$, see
Algebra, Lemma \ref{algebra-lemma-minimal-prime-reduced-ring}.
Thus $\text{d}_\eta$ is identified with the map
$$
\text{d} :
H^n(X_\eta, \Omega^{n - 1}_{X_\eta/\kappa(\eta)})
\longrightarrow
H^n(X_\eta, \Omega^n_{X_\eta/\kappa(\eta)})
$$
which is zero by Lemma \ref{lemma-top-part-degenerates}.
\end{proof}

\begin{proof}[Proof in the general case]
Observe that the question is flat local on $S$: if $S' \to S$ is a surjective
flat morphism of schemes and the map is zero after pullback to $S'$,
then the map is zero. Also, formation of the map commutes with base change
by flat morphisms by flat base change (Cohomology of Schemes, Lemma
\ref{coherent-lemma-flat-base-change-cohomology}).

\medskip\noindent
Consider the Stein factorization $X \to S' \to S$ as in
More on Morphisms, Theorem
\ref{more-morphisms-theorem-stein-factorization-general}.
By Lemma \ref{lemma-relative-bottom-part-degenerates} the morphism
$\pi : S' \to S$ is finite \'etale.
The morphism $f : X \to S'$ is proper (by the theorem),
smooth (by More on Morphisms, Lemma
\ref{more-morphisms-lemma-smooth-etale-permanence}) with geometrically
connected fibres by the theorem on Stein factorization.
In the proof of Lemma \ref{lemma-relative-bottom-part-degenerates}
we saw that $\Omega_{X/S} = \Omega_{X/S'}$ because $S' \to S$ is \'etale.
Hence $\Omega^\bullet_{X/S} = \Omega^\bullet_{X/S'}$.
We have
$$
R^qf_*\Omega^p_{X/S} = \pi_*R^qf'_*\Omega^p_{X/S'}
$$
for all $p, q$ by the Leray spectral sequence
(Cohomology, Lemma \ref{cohomology-lemma-relative-Leray}),
the fact that $\pi$ is finite hence affine, and
Cohomology of Schemes, Lemma \ref{coherent-lemma-relative-affine-vanishing}
(of course we also use that $R^qf'_*\Omega^p_{X'/S}$ is
quasi-coherent).
Thus the map of the lemma is $\pi_*$ applied to
$\text{d} : R^nf'_*\Omega^{n - 1}_{X/S'} \to R^nf'_*\Omega^n_{X/S'}$.
In other words, in order to prove the lemma we may replace
$f : X \to S$ by $f' : X \to S'$ to reduce to the case discussed
in the next pargraph.

\medskip\noindent
Assume $f$ has geometrically connected fibres and
$f_*\mathcal{O}_X = \mathcal{O}_S$.
For every $s \in S$ we can choose an \'etale neighbourhood
$(S', s') \to (S, s)$ such that the base change $X' \to S'$ of $S$
has a section. See More on Morphisms, Lemma
\ref{more-morphisms-lemma-etale-nbhd-dominates-smooth}.
By the initial remarks of the proof this reduces us to the case
discussed in the next paragraph.

\medskip\noindent
Assume $f$ has geometrically connected fibres,
$f_*\mathcal{O}_X = \mathcal{O}_S$, and we have
a section $s : S \to X$ of $f$. We may and do assume $S = \Spec(A)$
is affine. The map
$s^* : R\Gamma(X, \mathcal{O}_X) \to R\Gamma(S, \mathcal{O}_S) = A$
is a splitting of the map $A \to R\Gamma(X, \mathcal{O}_X)$. Thus we can write
$$
R\Gamma(X, \mathcal{O}_X) = A \oplus P
$$
where $P$ is the ``kernel'' of $s^*$. By
Lemma \ref{lemma-relative-bottom-part-degenerates} the object $P$
of $D(A)$ is perfect of tor amplitude in $[1, n]$. As in the proof
of Lemma \ref{lemma-relative-duality-hodge} we see that
$H^n(X, \Omega^n_{X/S})$ is a locally free $A$-module of rank $1$
(and in fact dual to $A$ so free of rank $1$ -- we will soon choose
a generator but we don't want to check it is the same generator
nor will it be necessary to do so).

\medskip\noindent
Denote $Z \subset X$ the image of $s$ which is a closed subscheme of $X$ by
Schemes, Lemma \ref{schemes-lemma-section-immersion}.
Observe that $Z \to X$ is a regular (and a fortiori Koszul regular by
Divisors, Lemma \ref{divisors-lemma-regular-quasi-regular-immersion})
closed immersion by
Divisors, Lemma \ref{divisors-lemma-section-smooth-regular-immersion}.
Of course $Z \to X$ has codimension $n$. Thus by
Remark \ref{remark-how-to-use}
we can consider the map
$$
\gamma^{0, 0} : H^0(Z, \Omega^0_{Z/S}) \longrightarrow H^n(X, \Omega^n_{X/S})
$$
and we set $\xi = \gamma^{0, 0}(1) \in H^n(X, \Omega^n_{X/S})$.

\medskip\noindent
We claim $\xi$ is a basis element. Namely, since we have base change in
top degree (see for example Limits, Lemma
\ref{limits-lemma-higher-direct-images-zero-above-dimension-fibre})
we see that
$H^n(X, \Omega^n_{X/S}) \otimes_A k = H^n(X_k, \Omega^n_{X_k/k})$
for any ring map $A \to k$. By the compatibility of
the construction of $\xi$ with base change,
see Lemma \ref{lemma-gysin-transverse-global},
we see that the image of $\xi$ in $H^n(X_k, \Omega^n_{X_k/k})$
is nonzero by Lemma \ref{lemma-class-of-a-point} if $k$ is a field.
Thus $\xi$ is a nowhere vanishing section of an invertible module
and hence a generator.

\medskip\noindent
Let $\theta \in H^n(X, \Omega^{n - 1}_{X/S})$. We have to show that
$\text{d}(\theta)$ is zero in $H^n(X, \Omega^n_{X/S})$.
We may write $\text{d}(\theta) = a \xi$ for some $a \in A$
as $\xi$ is a basis element. Then we have to show $a = 0$.

\medskip\noindent
Consider the closed immersion
$$
\Delta : X \to X \times_S X
$$
This is also a section of a smooth morphism (namely either projection)
and hence a regular and Koszul immersion of codimension $n$ as well.
Thus we can consider the maps
$$
\gamma^{p, q} :
H^q(X, \Omega^p_{X/S})
\longrightarrow
H^{q + n}(X \times_S X, \Omega^{p + n}_{X \times_S X/S})
$$
of Remark \ref{remark-how-to-use}. Consider the image
$$
\gamma^{n - 1, n}(\theta) \in
H^{2n}(X \times_S X, \Omega^{2n - 1}_{X \times_S X})
$$
By Lemma \ref{lemma-de-rham-complex-product} we have
$$
\Omega^{2n - 1}_{X \times_S X} =
\Omega^{n - 1}_{X/S} \boxtimes \Omega^n_{X/S} \oplus
\Omega^n_{X/S} \boxtimes \Omega^{n - 1}_{X/S}
$$
By the K\"unneth formula (either
Derived Categories of Schemes, Lemma \ref{perfect-lemma-kunneth} or
Derived Categories of Schemes, Lemma \ref{perfect-lemma-kunneth-single-sheaf})
we see that
$$
H^{2n}(X \times_S X, \Omega^{n - 1}_{X/S} \boxtimes \Omega^n_{X/S}) =
H^n(X, \Omega^{n - 1}_{X/S}) \otimes_A H^n(X, \Omega^n_{X/S})
$$
and
$$
H^{2n}(X \times_S X, \Omega^n_{X/S} \boxtimes \Omega^{n - 1}_{X/S}) =
H^n(X, \Omega^n_{X/S}) \otimes_A H^n(X, \Omega^{n - 1}_{X/S})
$$
Namely, since we are looking in top degree there no higher tor groups
that intervene. Combined with the fact that $\xi$ is a generator this means
we can write
$$
\gamma^{n - 1, n}(\theta) = \theta_1 \otimes \xi + \xi \otimes \theta_2
$$
with $\theta_1, \theta_2 \in H^n(X, \Omega^{n - 1}_{X/S})$.
Arguing in exactly the same manner we can write
$$
\gamma^{n, n}(\xi) = b \xi \otimes \xi
$$
in
$H^{2n}(X \times_S X, \Omega^{2n}_{X \times_S X/S}) =
H^n(X, \Omega^n_{X/S}) \otimes_A H^n(X, \Omega^n_{X/S})$
for some $b \in H^0(S, \mathcal{O}_S)$.

\medskip\noindent
{\bf Claim:} $\theta_1 = \theta$, $\theta_2 = \theta$, and $b = 1$.
Let us show that the claim implies the desired result $a = 0$.
Namely, by Lemma \ref{lemma-gysin-differential-hodge}
we have
$$
\gamma^{n, n}(\text{d}(\theta)) = \text{d}(\gamma^{n - 1, n}(\theta))
$$
By our choices above this gives
$$
a \xi \otimes \xi =
\gamma^{n, n}(a\xi) =
\text{d}(\theta \otimes \xi + \xi \otimes \theta) =
a \xi \otimes \xi + (-1)^n a \xi \otimes \xi
$$
The right most equality comes from the fact that the map
$\text{d} : \Omega^{2n - 1}_{X \otimes_S X/S} \to \Omega^{2n}_{X \times_S X/S}$
by Lemma \ref{lemma-de-rham-complex-product}
is the sum of the differential
$\text{d} \boxtimes 1 : \Omega^{n - 1}_{X/S} \boxtimes \Omega^n_{X/S}
\to \Omega^n_{X/S} \boxtimes \Omega^n_{X/S}$
and the differential
$(-1)^n 1 \boxtimes \text{d} : \Omega^n_{X/S} \boxtimes \Omega^{n - 1}_{X/S}
\to \Omega^n_{X/S} \boxtimes \Omega^n_{X/S}$. Please see discussion in
Section \ref{section-kunneth} and
Derived Categories of Schemes, Section
\ref{perfect-section-kunneth-complexes} for more information.
Since $\xi \otimes \xi$ is a basis for the rank $1$ free $A$-module
$H^n(X, \Omega^n_{X/S}) \otimes_A H^n(X, \Omega^n_{X/S})$
we conclude
$$
a = a + (-1)^n a \Rightarrow a = 0
$$
as desired.

\medskip\noindent
In the rest of the proof we prove the claim above. Let us denote
$\eta = \gamma^{0, 0}(1) \in H^n(X \times_S X, \Omega^n_{X \times_S X/S})$.
Since $\Omega^n_{X \times_S X/S} =
\bigoplus_{p + p' = n} \Omega^p_{X/S} \boxtimes \Omega^{p'}_{X/S}$
we may write
$$
\eta = \eta_0 + \eta_1 + \ldots + \eta_n
$$
where $\eta_p$ is in
$H^n(X \times_S X, \Omega^p_{X/S} \boxtimes \Omega^{n - p}_{X/S})$.
For $p = 0$ we can write
\begin{align*}
H^n(X \times_S X, \mathcal{O}_X \boxtimes \Omega^n_{X/S})
& =
H^n(R\Gamma(X, \mathcal{O}_X) \otimes_A^\mathbf{L}
R\Gamma(X, \Omega^n_{X/S})) \\
& =
A \otimes_A H^n(X, \Omega^n_{X/S}) \oplus
H^n(P \otimes_A^\mathbf{L} R\Gamma(X, \Omega^n_{X/S}))
\end{align*}
by our previously given decomposition $R\Gamma(X, \mathcal{O}_X) = A \oplus P$.
Consider the morphism $(s, \text{id}) : X \to X \times_S X$.
Then $(s, \text{id})^{-1}(\Delta) = Z$ scheme theoretically.
Hence we see that $(s, \text{id})^*\eta = \xi$ by
Lemma \ref{lemma-gysin-transverse-global}. This means that
$$
\xi = (s, \text{id})^*\eta = (s^* \otimes \text{id})(\eta_0)
$$
This means exactly that the first component of $\eta_0$
in the direct sum decomposition above is $\xi$. In other words, we can write
$$
\eta_0 = 1 \otimes \xi + \eta'_0
$$
with $\eta'_0 \in H^n(P \otimes_A^\mathbf{L} R\Gamma(X, \Omega^n_{X/S}))$.
In exactly the same manner for $p = n$ we can write
\begin{align*}
H^n(X \times_S X, \Omega^n_{X/S} \boxtimes \mathcal{O}_X)
& =
H^n(R\Gamma(X, \Omega^n_{X/S}) \otimes_A^\mathbf{L}
R\Gamma(X, \mathcal{O}_X)) \\
& =
H^n(X, \Omega^n_{X/S}) \otimes_A A \oplus
H^n(R\Gamma(X, \Omega^n_{X/S}) \otimes_A^\mathbf{L} P)
\end{align*}
and we can write
$$
\eta_n = \xi \otimes 1 + \eta'_n
$$
with $\eta'_n \in H^n(R\Gamma(X, \Omega^n_{X/S}) \otimes_A^\mathbf{L} P)$.

\medskip\noindent
Observe that $\text{pr}_1^*\theta = \theta \otimes 1$
and $\text{pr}_2^*\theta = 1 \otimes \theta$ are
Hodge cohomology classes on
$X \times_S X$ which pull back to $\theta$ by $\Delta$.
Hence by Lemma \ref{lemma-gysin-projection-global} we have
$$
\theta_1 \otimes \xi + \xi \otimes \theta_2 =
\gamma^{n - 1, n}(\theta) =
(\theta \otimes 1) \cup \eta =
(1 \otimes \theta) \cup \eta
$$
in the Hodge cohomology ring of $X \times_S X$ over $S$.
In terms of the direct sum decomposition on the modules
of differentials of $X \times_S X/S$ we obtain
$$
\theta_1 \otimes \xi =
(\theta \otimes 1) \cup \eta_0
\quad\text{and}\quad
\xi \otimes \theta_2 =
(1 \otimes \theta) \cup \eta_n
$$
Looking at the formula $\eta_0 = 1 \otimes \xi + \eta'_0$ we found above,
we see that to show that $\theta_1 = \theta$ it suffices to prove that
$$
(\theta \otimes 1) \cup \eta'_0 = 0
$$
To do this, observe that cupping with $\theta \otimes 1$ is given
by the action on cohomology of the map
$$
(P \otimes_A^\mathbf{L} R\Gamma(X, \Omega^n_{X/S}))[-n]
\xrightarrow{\theta \otimes 1}
R\Gamma(X, \Omega^{n - 1}_{X/S}) \otimes_A^\mathbf{L}
R\Gamma(X, \Omega^n_{X/S})
$$
in the derived category, see Cohomology, Remark
\ref{cohomology-remark-cup-with-element-map-total-cohomology}.
This map is the derived tensor product of the two maps
$$
\theta : P[-n] \to R\Gamma(X, \Omega^{n - 1}_{X/S})
\quad\text{and}\quad
1 : R\Gamma(X, \Omega^n_{X/S}) \to R\Gamma(X, \Omega^n_{X/S})
$$
by Derived Categories of Schemes, Remark
\ref{perfect-remark-annoying-compatibility}.
However, the first of these is zero in $D(A)$ because it is a map from
a perfect complex of tor amplitude in $[n + 1, 2n]$ to a complex
with cohomology only in degrees $0, 1, \ldots, n$, see
More on Algebra, Lemma \ref{more-algebra-lemma-splitting-unique}.
A similar argument works to show the vanishing of
$(1 \otimes \theta) \cup \eta'_n$. Finally, 
in exactly the same manner we obtain
$$
b \xi \otimes \xi = \gamma^{n, n}(\xi) = (\xi \otimes 1) \cup \eta_0
$$
and we conclude as before by showing that
$(\xi \otimes 1) \cup \eta'_0 = 0$ in the same manner as above.
This finishes the proof.
\end{proof}

\begin{proposition}
\label{proposition-relative-poincare-duality}
Let $S$ be a quasi-compact and quasi-separated scheme. Let $f : X \to S$
be a proper smooth morphism of schemes all of whose fibres are nonempty
and equidimensional of dimension $n$. There exists an
$\mathcal{O}_S$-module map
$$
t : R^{2n}f_*\Omega^\bullet_{X/S} \longrightarrow \mathcal{O}_S
$$
unique up to precomposing by multiplication by a unit of
$H^0(X, \mathcal{O}_X)$ with the following property: the pairing
$$
Rf_*\Omega^\bullet_{X/S}
\otimes_{\mathcal{O}_S}^\mathbf{L}
Rf_*\Omega^\bullet_{X/S}[2n]
\longrightarrow
\mathcal{O}_S, \quad
(\xi, \xi') \longmapsto t(\xi \cup \xi')
$$
is a perfect pairing of perfect complexes on $S$.
\end{proposition}

\begin{proof}
The proof is exactly the same as the proof of
Proposition \ref{proposition-poincare-duality}.

\medskip\noindent
By the relative Hodge-to-de Rham spectral sequence
$$
E_1^{p, q} = R^qf_*\Omega^p_{X/S} \Rightarrow R^{p + q}f_*\Omega^\bullet_{X/S}
$$
(Section \ref{section-hodge-to-de-rham}), the vanishing
of $\Omega^i_{X/S}$ for $i > n$, the vanishing in for example Limits, Lemma
\ref{limits-lemma-higher-direct-images-zero-above-dimension-fibre}
and the results of Lemmas \ref{lemma-relative-bottom-part-degenerates} and
\ref{lemma-relative-top-part-degenerates}
we see that $R^0f_*\Omega_{X/S} = R^0f_*\mathcal{O}_X$
and $R^nf_*\Omega^n_{X/S} = R^{2n}f_*\Omega^\bullet_{X/S}$.
More precisely, these identifications come from the maps
of complexes
$$
\Omega^\bullet_{X/S} \to \mathcal{O}_X[0]
\quad\text{and}\quad
\Omega^n_{X/S}[-n] \to \Omega^\bullet_{X/S}
$$
Let us choose $t : R^{2n}f_*\Omega_{X/S} \to \mathcal{O}_S$
which via this identification corresponds to a $t$ as in
Lemma \ref{lemma-relative-duality-hodge}.

\medskip\noindent
Let us abbreviate $\Omega^\bullet = \Omega^\bullet_{X/S}$.
Consider the map (\ref{equation-wedge}) which in our situation reads
$$
\wedge :
\text{Tot}(\Omega^\bullet \otimes_{f^{-1}\mathcal{O}_S} \Omega^\bullet)
\longrightarrow
\Omega^\bullet
$$
For every integer $p = 0, 1, \ldots, n$ this map annihilates the subcomplex
$\text{Tot}(\sigma_{> p} \Omega^\bullet \otimes_{f^{-1}\mathcal{O}_S}
\sigma_{\geq n - p} \Omega^\bullet)$ for degree reasons.
Hence we find that the restriction of $\wedge$ to the subcomplex
$\text{Tot}(\Omega^\bullet \otimes_{f^{-1}\mathcal{O}_S}
\geq_{n - p}\Omega^\bullet)$ factors through a map of complexes
$$
\gamma_p :
\text{Tot}(\sigma_{\leq p} \Omega^\bullet \otimes_{f^{-1}\mathcal{O}_S}
\sigma_{\geq n - p} \Omega^\bullet)
\longrightarrow
\Omega^\bullet
$$
Using the same procedure as in Section \ref{section-cup-product} we obtain
relative cup products
$$
Rf_*\sigma_{\leq p} \Omega^\bullet
\otimes_{\mathcal{O}_S}^\mathbf{L}
Rf_*\sigma_{\geq n - p}\Omega^\bullet
\longrightarrow
Rf_*\Omega^\bullet
$$
We will prove by induction on $p$ that these cup products via $t$
induce perfect pairings between $Rf_*\sigma_{\leq p} \Omega^\bullet$
and $Rf_*\sigma_{\geq n - p}\Omega^\bullet[2n]$. For $p = n$
this is the assertion of the proposition.

\medskip\noindent
The base case is $p = 0$. In this case we have
$$
Rf_*\sigma_{\leq p}\Omega^\bullet = Rf_*\mathcal{O}_X
\quad\text{and}\quad
Rf_*\sigma_{\geq n - p}\Omega^\bullet[2n] = Rf_*(\Omega^n[-n])[2n] =
Rf_*\Omega^n[n]
$$
In this case we simply obtain the pairing
between $Rf_*\mathcal{O}_X$ and $Rf_*\Omega^n[n]$ of
Lemma \ref{lemma-relative-duality-hodge} and the result is true.

\medskip\noindent
Induction step. Say we know the result is true for $p$. Then
we consider the distinguished triangle
$$
\Omega^{p + 1}[-p - 1] \to
\sigma_{\leq p + 1}\Omega^\bullet \to
\sigma_{\leq p}\Omega^\bullet \to
\Omega^{p + 1}[-p]
$$
and the distinguished triangle
$$
\sigma_{\geq n - p}\Omega^\bullet \to
\sigma_{\geq n - p - 1}\Omega^\bullet \to
\Omega^{n - p - 1}[-n + p + 1] \to
(\sigma_{\geq n - p}\Omega^\bullet)[1]
$$
Observe that both are distinguished triangles in the homotopy category
of complexes of sheaves of $f^{-1}\mathcal{O}_S$-modules; in particular the
maps $\sigma_{\leq p}\Omega^\bullet \to \Omega^{p + 1}[-p]$ and
$\Omega^{n - p - 1}[-d + p + 1] \to (\sigma_{\geq n - p}\Omega^\bullet)[1]$
are given by actual maps of complexes, namely using the differential
$\Omega^p \to \Omega^{p + 1}$ and the differential
$\Omega^{n - p - 1} \to \Omega^{n - p}$.
Consider the distinguished triangles associated gotten from these
distinguished triangles by applying $Rf_*$
$$
\xymatrix{
Rf_*\sigma_{\leq p}\Omega^\bullet \ar[d]_a \\
Rf_*\Omega^{p + 1}[-p - 1] \ar[d]_b \\
Rf_*\sigma_{\leq p + 1}\Omega^\bullet \ar[d]_c \\
Rf_*\sigma_{\leq p}\Omega^\bullet \ar[d]_d \\
Rf_*\Omega^{p + 1}[-p - 1]
}
\quad\quad
\xymatrix{
Rf_*\sigma_{\geq n - p}\Omega^\bullet \\
Rf_*\Omega^{n - p - 1}[-n + p + 1] \ar[u]_{a'} \\
Rf_*\sigma_{\geq n - p - 1}\Omega^\bullet \ar[u]_{b'} \\
Rf_*\sigma_{\geq n - p}\Omega^\bullet \ar[u]_{c'} \\
Rf_*\Omega^{n - p - 1}[-n + p + 1] \ar[u]_{d'}
}
$$
We will show below that the pairs $(a, a')$, $(b, b')$, $(c, c')$, and
$(d, d')$ are compatible with the given pairings. This means we obtain a
map from the distinguished triangle on the left to the distinguished triangle
obtained by applying $R\SheafHom(-, \mathcal{O}_S)$ to the distinguished
triangle on the right. By induction and Lemma \ref{lemma-duality-hodge}
we know that the pairings constructed above between the
complexes on the first, second, fourth, and fifth
rows are perfect, i.e., determine isomorphisms after taking duals.
By Derived Categories, Lemma \ref{derived-lemma-third-isomorphism-triangle}
we conclude the pairing between the complexes in the middle row
is perfect as desired.

\medskip\noindent
Let $e : K \to K'$ and $e' : M' \to M$ be maps of objects
of $D(\mathcal{O}_S)$ and let
$K \otimes_{\mathcal{O}_S}^\mathbf{L} M \to \mathcal{O}_S$ and
$K' \otimes_{\mathcal{O}_S}^\mathbf{L} M' \to \mathcal{O}_S$
be pairings. Then we say these pairings are compatible if the
diagram
$$
\xymatrix{
K' \otimes_{\mathcal{O}_S}^\mathbf{L} M' \ar[d] &
K \otimes_{\mathcal{O}_S}^\mathbf{L} M'
\ar[l]^{e \otimes 1} \ar[d]^{1 \otimes e'} \\
\mathcal{O}_S &
K \otimes_{\mathcal{O}_S}^\mathbf{L} M \ar[l]
}
$$
commutes. This indeed means that the diagram
$$
\xymatrix{
K \ar[r] \ar[d]_e & R\SheafHom(M, \mathcal{O}_S) \ar[d]^{R\SheafHom(e', -)} \\
K' \ar[r] & R\SheafHom(M', \mathcal{O}_S)
}
$$
commutes and hence is sufficient for our purposes.

\medskip\noindent
Let us prove this for the pair $(c, c')$. Here we observe simply
that we have a commutative diagram
$$
\xymatrix{
\text{Tot}(\sigma_{\leq p} \Omega^\bullet \otimes_{f^{-1}\mathcal{O}_S}
\sigma_{\geq n - p} \Omega^\bullet) \ar[d]_{\gamma_p} &
\text{Tot}(\sigma_{\leq p + 1} \Omega^\bullet \otimes_{f^{-1}\mathcal{O}_S}
\sigma_{\geq n - p} \Omega^\bullet) \ar[l] \ar[d] \\
\Omega^\bullet &
\text{Tot}(\sigma_{\leq p + 1} \Omega^\bullet \otimes_{f^{-1}\mathcal{O}_S}
\sigma_{\geq n - p - 1} \Omega^\bullet) \ar[l]_-{\gamma_{p + 1}}
}
$$
By functoriality of the cup product we obtain commutativity of the
desired diagram.

\medskip\noindent
Similarly for the pair $(b, b')$ we use the commutative diagram
$$
\xymatrix{
\text{Tot}(\sigma_{\leq p + 1} \Omega^\bullet \otimes_{f^{-1}\mathcal{O}_S}
\sigma_{\geq n - p - 1} \Omega^\bullet) \ar[d]_{\gamma_{p + 1}} &
\text{Tot}(\Omega^{p + 1}[-p - 1] \otimes_{f^{-1}\mathcal{O}_S}
\sigma_{\geq n - p - 1} \Omega^\bullet) \ar[l] \ar[d] \\
\Omega^\bullet &
\Omega^{p + 1}[-p - 1]
\otimes_{f^{-1}\mathcal{O}_S}
\Omega^{n - p - 1}[-n + p + 1] \ar[l]_-\wedge
}
$$

\medskip\noindent
For the pairs $(d, d')$ and $(a, a')$ we use the commutative diagram
$$
\xymatrix{
\Omega^{p + 1}[-p] \otimes_{f^{-1}\mathcal{O}_S}
\Omega^{n - p - 1}[-n + p] \ar[d] &
\text{Tot}(\sigma_{\leq p}\Omega^\bullet \otimes_{f^{-1}\mathcal{O}_S}
\Omega^{n - p - 1}[-n + p]) \ar[l] \ar[d] \\
\Omega^\bullet &
\text{Tot}(\sigma_{\leq p}\Omega^\bullet \otimes_{f^{-1}\mathcal{O}_S}
\sigma_{\geq n - p}\Omega^\bullet) \ar[l]
}
$$

\medskip\noindent
We omit the argument showing the uniqueness of $t$ up to
precomposing by multiplication by a unit in $H^0(X, \mathcal{O}_X)$.
\end{proof}

\bibliographystyle{plain}
\bibliography{my}
\end{document}
